% Partie : logique du premier ordre
% Chapitre : complétude
% Section : compacité

\documentclass[../../../include/open-logic-section]{subfiles}

\begin{document}

\iftag{FOL}
      {\olfileid[fr]{fol}{com}{com}}
      {\olfileid[fr]{pl}{com}{com}}

\olsection{Le théorème de compacité}

Le théorème de compacité est une conséquence importante du théorème de
complétude. Il affirme que si chaque sous-ensemble \emph{fini} d'un
ensemble \iftag{FOL}{d'!!{sentence}s}{de !!{formula}s} est satisfaisable, l'ensemble entier l'est
aussi, même s'il est infini. Cela n'a rien d'évident. À première vue,
on pourrait imaginer qu'un ensemble infini \iftag{FOL}{d'!!{sentence}s}{de !!{formula}s} soit
contradictoire uniquement en raison de son infinité, alors que chacun
de ses sous-ensembles finis serait satisfaisable. Le théorème exclut
cette possibilité : aucun ensemble infini non satisfaisable \iftag{FOL}{d'!!{sentence}s}{de !!{formula}s} ne possède cette propriété. Comme pour la complétude,
il existe aussi une formulation en termes de conséquence logique : si
un ensemble infini \iftag{FOL}{d'!!{sentence}s}{de !!{formula}s} entraîne \iftag{FOL}{un !!{sentence}}{une !!{formula}}, un de ses
sous-ensembles finis l'entraîne déjà.

\begin{defn}
  Un ensemble~$\Gamma$ \iftag{FOL}{d'!!{sentence}s}{de !!{formula}s}
  est \emph{finiment satisfaisable} si et seulement si tout
  sous-ensemble fini $\Gamma_0\subseteq\Gamma$ est satisfaisable.
\end{defn}

\begin{thm}[Théorème de compacité]
\ollabel{thm:compactness}
Soient $\Gamma$ un ensemble \iftag{FOL}{d'!!{sentence}s}{de !!{formula}s} et $!A$
\iftag{FOL}{un !!{sentence}}{une !!{formula}} du même langage. Alors :
\begin{enumerate}
  \item $\Gamma\Entails !A$ si et seulement s'il existe un
    sous-ensemble fini $\Gamma_0\subseteq\Gamma$ tel que
    $\Gamma_0\Entails !A$.
  \item $\Gamma$ est satisfaisable si et seulement s'il est finiment
    satisfaisable.
\end{enumerate}
\end{thm}

\begin{proof}
Démontrons (2). Si $\Gamma$ est satisfaisable, il existe
\iftag{FOL}{!!a{structure}~$\Struct{M}$}{!!a{valuation}~$\pAssign{v}$}
telle que $\iftag{FOL}{\Sat{M}}{\pSat{v}}{!A}$ pour \iftag{FOL}{tout !!{sentence}}{toute !!{formula}}
$!A\in\Gamma$. Cette $\iftag{FOL}{\Struct{M}}{\pAssign{v}}$
satisfait évidemment aussi chaque sous-ensemble fini de~$\Gamma$ ;
donc $\Gamma$ est finiment satisfaisable.

Supposons maintenant $\Gamma$ finiment satisfaisable. Alors chaque
sous-ensemble fini $\Gamma_0\subseteq\Gamma$ est satisfaisable.
Par correction
(\iftag{FOL}{%
  \tagrefs{prfAX/{fol:axd:sou:cor:consistency-soundness},
    prfSC/{fol:seq:sou:cor:consistency-soundness},
    prfND/{fol:ntd:sou:cor:consistency-soundness},
    prfTab/{fol:tab:sou:cor:consistency-soundness}}}{%
  \tagrefs{prfAX/{pl:axd:sou:cor:consistency-soundness},
    prfSC/{pl:seq:sou:cor:consistency-soundness},
    prfND/{pl:ntd:sou:cor:consistency-soundness},
    prfTab/{pl:tab:sou:cor:consistency-soundness}}}),
chaque sous-ensemble fini est cohérent. Il s'ensuit que $\Gamma$
lui-même est cohérent, d'après les
\iftag{FOL}{%
  \tagrefs{prfAX/{fol:axd:ptn:prop:proves-compact},
    prfSC/{fol:seq:ptn:prop:proves-compact},
    prfND/{fol:ntd:ptn:prop:proves-compact},
    prfTab/{fol:tab:ptn:prop:proves-compact}}}{%
  \tagrefs{prfAX/{pl:axd:ptn:prop:proves-compact},
    prfSC/{pl:seq:ptn:prop:proves-compact},
    prfND/{pl:ntd:ptn:prop:proves-compact},
    prfTab/{pl:tab:ptn:prop:proves-compact}}}.
Par complétude (\olref[cth]{thm:completeness}), puisque $\Gamma$ est
cohérent, il est satisfaisable.
\end{proof}

\tagprob{FOL}
\begin{prob}
Démontrez (1) du~\olref[fol][com][com]{thm:compactness}.
\end{prob}
\tagendprob

\tagprob{notFOL}
\begin{prob}
Démontrez (1) du~\olref[pl][com][com]{thm:compactness}.
\end{prob}
\tagendprob

\iftag{FOL}{%
\begin{ex}
Dans tout modèle~$\Struct{M}$ d'une théorie~$\Gamma$, chaque terme
\emph{clos}~$t$ désigne bien sûr !!a{element} de~$\Domain{M}$.
Peut-on garantir que, réciproquement, chaque !!{element} de
$\Domain{M}$ est désigné par un terme clos ? Autrement dit,
existe-t-il des théories~$\Gamma$ dont tous les modèles sont
couverts ? Le théorème de compacité montre que ce n'est pas le cas
si~$\Gamma$ possède un modèle infini. Voici pourquoi. Soit
$\Struct{M}$ un modèle infini de~$\Gamma$, et soit $c$ !!a{constant}
étrangère au langage de~$\Gamma$. Soit $\Delta$ l'ensemble de tous
les énoncés $\eq/[c][t]$ où $t$ est un terme \emph{clos} du
langage~$\Lang{L}$ de~$\Gamma$, c'est-à-dire
  \[
  \Delta=\Setabs{\eq/[c][t]}{t\in\Trm[L]\text{ et }t\text{ est clos}}.
  \]
Un sous-ensemble fini de $\Gamma\cup\Delta$ s'écrit
$\Gamma'\cup\Delta'$, où $\Gamma'\subseteq\Gamma$ et
$\Delta'\subseteq\Delta$. Puisque $\Delta'$ est fini, il n'y figure
qu'un nombre fini de termes. Choisissons un
$a\in\Domain{M}$ qui n'est désigné par aucun d'eux, et soit
$\Struct{M'}$ la !!{structure} obtenue en interprétant en outre
$c$ par $\Assign{c}{M'}=a$. Comme
$a\ne\Value{t}{M}$ pour chaque terme~$t$ figurant dans~$\Delta'$,
on a $\Sat{M'}{\Delta'}$. De plus,
$\Sat{M}{\Gamma}$, $\Gamma'\subseteq\Gamma$, et $c$ n'apparaît pas
dans~$\Gamma$ ; donc $\Sat{M'}{\Gamma'}$. Ainsi,
$\Sat{M'}{\Gamma'\cup\Delta'}$ pour chaque sous-ensemble fini
$\Gamma'\cup\Delta'$ de~$\Gamma\cup\Delta$. Par compacité,
$\Gamma\cup\Delta$ est satisfaisable. Soit $\Struct{M^+}$ l'un de
ses modèles. Le réduit de~$\Struct{M^+}$ au langage~$\Lang L$ est
un modèle de~$\Gamma$ qui n'est pas couvert : l'élément
$\Assign{c}{M^+}$ de son domaine est distinct de la valeur de tout
terme clos~$t$ de~$\Lang L$.
\end{ex}

\begin{ex}
Considérons !!a{language}~$\Lang{L}$ contenant le
!!{predicate}~$<$, les !!{constant}s $\Obj{0}$ et $\Obj{1}$, et
les !!{function}s $+$, $\times$ et $-$. Soit $\Gamma$ l'ensemble
des !!{sentence}s de ce langage qui sont vrais dans la
!!{structure}~$\Struct{Q}$ de domaine~$\Rat$, avec les
interprétations usuelles. Ainsi, $\Gamma$ exprime tous les énoncés
de~$\Lang{L}$ vrais dans les rationnels. Dans~$\Rat$ (et même
dans~$\Real$), aucun nombre~$r$ n'est supérieur à~$0$ tout en étant
inférieur à~$1/k$ pour chaque $k\in\PosInt$. Un tel nombre, s'il
existait, serait un \emph{infinitésimal strictement positif}. Le théorème de compacité permet de montrer qu'il existe des
modèles de~$\Gamma$ contenant des infinitésimaux. Notre langage n'a
pas de !!{function} de division (la division par zéro n'est pas
définie, alors que les !!{function}s doivent être interprétés par
des fonctions totales). Nous pouvons néanmoins exprimer $r<1/k$ par
$r\cdot k<1$. Soit maintenant $c$ une nouvelle !!{constant}, et
posons
\[
\Delta=\{\Obj{0}<c\}
\cup\Setabs{c\times\num k<\Obj{1}}{k\in\PosInt}
\]
(où $\num{k}=(\Obj{1}+(\Obj{1}+\dots+(\Obj{1}+
\Obj{1})\dots))$ contient $k$ occurrences de~$\Obj{1}$).
Pour tout sous-ensemble fini~$\Delta_0$ de~$\Delta$, il existe un
entier~$K>0$ tel que chaque énoncé $c\times\num{k}<\Obj{1}$
appartenant à~$\Delta_0$ ait $k<K$. Si nous élargissons
$\Struct{Q}$ en~$\Struct{Q'}$ avec
$\Assign{c}{Q'}=1/K$, alors
$\Sat{Q'}{\Gamma_0\cup\Delta_0}$ pour tout sous-ensemble fini
$\Gamma_0\subseteq\Gamma$. Ainsi, $\Gamma\cup\Delta$ est finiment
satisfaisable (exercice : démontrez-le en détail). Par compacité,
$\Gamma\cup\Delta$ est satisfaisable. Tout modèle~$\Struct{S}$ de
$\Gamma\cup\Delta$ contient un infinitésimal, à savoir
$\Assign{c}{S}$.
\end{ex}
}{}

\tagprob{FOL}
\begin{prob}
Dans le modèle standard de l'arithmétique~$\Struct{N}$, aucun
!!{element}~$k\in\Domain{N}$ ne satisfait toutes les formules
$\num{n}<x$ (où $\num{n}$ vaut
$\Obj{0}^{\prime\dots\prime}$ avec $n$ occurrences de~$\prime$).
Utilisez le théorème de compacité pour montrer que les !!{sentence}s
du langage de l'arithmétique vrais dans le modèle standard
$\Struct{N}$ sont aussi vrais dans !!a{structure}~$\Struct{N'}$
qui contient !!a{element} satisfaisant chacune des formules
$\num{n}<x$.
\end{prob}
\tagendprob

\iftag{FOL}{
\begin{ex}
Nous savons qu'avec l'!!{identity}, la logique du premier ordre peut
imposer une taille minimale au !!{domain}. L'énoncé~$!A_{\ge n}$,
qui dit « il y a au moins $n$ objets distincts », est vrai
exactement dans les structures dont le domaine possède au moins~$n$
éléments.
Si nous posons
  \[
  \Delta=\Setabs{!A_{\ge n}}{n\ge1},
  \]
alors tout modèle de~$\Delta$ est infini. Ajouter~$\Delta$ à une
théorie permet donc de ne conserver que ses modèles infinis : les
modèles de $\Gamma\cup\Delta$ sont exactement les modèles infinis
de~$\Gamma$.

Une \emph{théorie}, c'est-à-dire un ensemble d'énoncés du premier
ordre, peut ainsi imposer l'infinitude. Le théorème de compacité
montre toutefois qu'aucun ensemble d'énoncés ne peut caractériser
exactement la finitude. Supposons en effet qu'un ensemble d'énoncés
$\Lambda$ soit satisfait dans toutes les !!{structure}s finies et
seulement dans celles-ci. Alors $\Delta\cup\Lambda$ est finiment
satisfaisable. Pourquoi ? Soit
$\Delta'\cup\Lambda'\subseteq\Delta\cup\Lambda$ un sous-ensemble
fini, avec $\Delta'\subseteq\Delta$ et
$\Lambda'\subseteq\Lambda$. Prenons pour~$n$ le plus grand entier
tel que $!A_{\ge n}\in\Delta'$ ; si $\Delta'$ est vide, prenons
$n=0$. Puisque $\Lambda$ est satisfait dans toutes les
structures finies, il existe un modèle~$\Struct{M}$ de~$\Lambda$
dont le domaine est fini et possède au moins~$n$ !!{element}s.
Alors $\Sat{M}{\Delta'\cup\Lambda'}$. Par compacité,
$\Delta\cup\Lambda$ possède un modèle infini, contrairement à
l'hypothèse selon laquelle $\Lambda$ n'est satisfait que dans des
!!{structure}s finies.
\end{ex}
}{}

\end{document}
